\documentclass[10pt,a4paper]{amsart}
\usepackage[margin=19mm]{geometry}
\usepackage{amsmath,amssymb,mathtools}
\usepackage[scr=rsfs,cal=euler]{mathalfa}
\usepackage{xfrac}
\usepackage[unicode,hidelinks,bookmarks=false]{hyperref}
\newcommand{\ie}{i.e.\ }
\newcommand{\cf}{cf.\ }
\newcommand{\resp}{respectively\ }
\newcommand{\Subsection}{\subsection}
\providecommand{\nobreakdash}{\nobreak\hbox{-}}
\setlength{\emergencystretch}{2em}
\multlinegap0pt
\usepackage{bm}
\usepackage{bbm}
\usepackage{dsfont}
\usepackage{xspace}
\usepackage{cancel}
\usepackage{mathtools}
\usepackage{amsthm}
\makeatletter
\@ifundefined{lemma}{\newtheorem{lemma}{Lemma}}{}
\makeatother
\renewcommand{\geq}{\geqslant}
\renewcommand{\leq}{\leqslant}
\title[Two dimensional covering systems and possible prime producin]{Two dimensional covering systems and possible prime producing $a^m-b^n$}
\author{Andrew Granville, Francesco Pappalardi}
\date{Source: arXiv:2601.10296v2 (CC BY 4.0)}
\begin{document}
\maketitle

\noindent\emph{Source and licence.} Two dimensional covering systems and possible prime producing $a^m-b^n$. Version arXiv:2601.10296v2 (CC BY 4.0), distributed under the Creative Commons Attribution 4.0 International licence, as recorded in the arXiv licence metadata for this exact version and shown on the abstract page https://arxiv.org/abs/2601.10296v2. Full rights notice: licenses/arxiv-2601.10296-CC-BY-4.0.txt.\par\smallskip

\noindent\textit{Editorial scope.} Counted unit: the Lemma whose source label is \texttt{lem: 1} in the article (source0.tex lines 256--259, source label \texttt{lem: 1}), delivered together with the article's own complete original proof of it (source0.tex lines 261--279, one \texttt{proof} environment of 19 lines). No mathematics is written, repaired or re-derived: statement and proof are the article's own text line for line. Required context retained uncounted: none. Every definition, display and label of those lines is retained unchanged. Omitted from this package and not redistributed: 1--255, 260--260, 280--910. Nothing inside a retained excerpt is omitted or altered: every retained body is delivered exactly as the article's lines are written, so no omission receipt of any kind accompanies this package. Citations: the retained excerpts carry no citation command at all, so no bibliography entry of the article is transcribed and no reference is redistributed. Reference adaptation: none was needed and none was made; every reference inside the delivered unit resolves inside it, so no unresolved reference or citation is produced. Printed slips kept literal: no printed word, symbol, hypothesis or number of the counted statement or of its proof was altered. Layout and class: the article's own class is \texttt{amsart} with packages including xy, geometry, amssymb, amsmath, amsthm, lineno, color; the wrapper is the lane's pinned \texttt{amsart} editorial wrapper at 10pt on A4, which restates the theorem-like environments the retained text needs and adds the article's own macro definitions that the retained text needs (\texttt{\textbackslash geq}, \texttt{\textbackslash leq}), with extra wrapper packages (bm, bbm, dsfont, xspace, cancel, mathtools). The delivered numbering is the unit's own. Assets: no retained span contains a figure environment, a graphics-inclusion command, a PDF-page inclusion, a table or a TikZ picture; no image file is copied into this package, the package declares no asset dependency and no image file is redistributed with it. Licence: version arXiv:2601.10296v2 of this article is distributed under the Creative Commons Attribution 4.0 International licence, as recorded in the arXiv licence metadata for this exact version and shown on the abstract page https://arxiv.org/abs/2601.10296v2. A full rights notice is delivered at licenses/arxiv-2601.10296-CC-BY-4.0.txt. Characters: every retained excerpt is pure ASCII, so the delivered engine, XeLaTeX with its default Latin Modern text font, reports no missing character.
\begin{lemma}  \label{lem: 1} (a) Every triple $(u,v,r)$ with $r\geq 1$  is equivalent to a semi-reduced triple. \\
% Let $u,v,r$ be given integers with $(u,v,r)=1$ and define $U=(u,r)$ so that $U$ divides $r$. There exists an integer $V$ with $(U,V)=1$ such that $S(u,v,r)=S(U,V,r)$.   In fact we select $V$ to be the smallest positive integer in a certain specified residue class mod $r/U$ for which $(V,U)=1$
(b)   If  $(u,v,r)$ and $(U,V,R)$ are semi-reduced then $S(u,v,r)=S(U,V,R)$ if and only if $u=U, r=R$ and $V\equiv v  \pmod {r/u}$ with $(V,u)=1$.
  \end{lemma}

 \begin{proof}  (a):  We begin by replacing $u$ and $v$ by their least positive residues $\pmod r$.

 We may assume that gcd$(u,v,r)=1$ for if $g=$gcd$(u,v,r)$ and $u=gU, v=gV, r=gR$ then $mv\equiv nu \pmod r$ if and only if $mV\equiv nU \pmod R$ and so $S(u,v,r)=S(U,V,R)$.
 
 
 We now show that we may assume that $u$ divides $r$: If not, let
  $U=(u,r)$ and write $u=Uh, r=UR$ with $(h,R)=1$. Now $(U,v)=(u,v,r)=1$ and so if $mv\equiv nu \pmod r$ then $U$ divides $mv$ and so $U$ divides $m$. Writing $m=MU$, we have $mv\equiv nu \pmod r$ if and only if 
 $Mv\equiv n h \pmod { R}$  which holds  if and only if  $n \equiv M vk \pmod { R}$ where
$k$ is the inverse of $h \pmod R$. % Let $U'$ be the product of the primes dividing $r$ that do not divide $R$ (so that $U'$ divides $U$), and
Select $V$ to be a positive integer $\leq   RU=r$ such that $V\equiv vk \pmod R$ and $(V,U)=1$. Therefore $n \equiv M vk \pmod { R}$ if and only if $n \equiv MV \pmod { R}$ and then, multiplying through by $U,$ this holds if and only if $nU \equiv m V \pmod r$.  Thus $S(u,v,r)=S(U,V,r)$ and  $ U=(u,r)$ divides $r$ with $(U,V)=1$ and  $1\leq V\leq r$.  

We now deduce that $(u,v)=(u,v,r)=1$ as $(u,r)=u$.  Moreover given any  triple $(u,v,r)$ for which $u$ divides $r$ and  $(u,v)=1$, we replace 
$v$ by $V$ the smallest positive residue of $v \pmod r$ and we obtain a semi-reduced triple $(u,V,r)$ with $S(u,v,r)=S(u,V,r)$.
\medskip

(b): We are assuming that $mv\equiv nu \pmod r$ if and only if $mV\equiv nU \pmod R$.
Now $u$ divides $r$ so $u$ divides $mv$ and therefore $m$ as $(u,v)=1$. Therefore $(m,n)=(u,v)$ gives the smallest $m$-value in a solution, and so $u=U$. More generally if $m=u$ then we see that the arithmetic progressions $v  \pmod {r/u}$ and $V \pmod {R/u}$
contain the same integers so $r/u=R/U$ and thus $r=R$ and $V\equiv v \pmod {r/u}$.
   \end{proof}

\end{document}
